\begin{abstract}
Public basis labels in BB84 retain reliability information hidden by aggregate quantum bit
error rate. We ask when that information survives estimation and coding costs. Complete
circuit paths explain the asymmetry, including the receiver rotation, while independent
channel noise progressively removes it. A nested LDPC family, pooled-centred shrinkage,
separate calibration data, and 24,000 main holdout frames establish a low-failure operating
regime. When every preparation path is noisy, a matched-code comparison at 10\% target QBER with a stronger
decoder reduces frame errors from 1.40\% to 0.65\% (ratio 0.46; 95\% interval 0.26--0.69).
Balanced paths show no corresponding benefit, and conservative calibration can increase
disclosure despite useful basis information. A complementary equal-budget study compares
noise scaling with independent calibration and a published interception-density estimator.
Wider, analytically selected noise settings improve extrapolation precision; drift and
unscaled channel disturbance expose different limits on interpretation. Together, the
experiments connect circuit exposure, finite-sample reliability, and acquisition design,
providing reproducible engineering guidance for gate-based BB84. Reconciliation yield and
model-based disturbance estimation remain distinct from a secure-key rate.
\end{abstract}
@METADATA@
\maketitle
\raggedbottom

% The narrative follows the information budget: public labels first, extra
% acquisitions second. Numerical claims are generated from one analysis.
\section{Introduction}
Two BB84 links can have the same overall quantum bit error rate (QBER) and present different
decoding problems. Errors may be spread evenly across their sifted bits or concentrated in
one basis. Public basis announcements retain this distinction. A reconciliation decoder can
use it through position-dependent probabilities, provided those probabilities can be learned
accurately enough from the bits sacrificed for estimation.

Gate-based implementations make the source of such differences explicit. Alternative
preparation sequences accumulate different noise exposure, and a receiver's basis rotation
can restore an asymmetry that preparation balancing appears to remove. This offers a useful
controlled setting for a broader engineering question: when does additional reliability
information reduce the cost of reconciling a frame, once code quality and estimation error
are held accountable?

The answer requires more than comparing failure ratios at a poorly chosen code rate. We
first identify the complete noisy paths, then calibrate a structured code family for a
low-FER operating regime using data separate from the final evaluation. We also centre
class-specific probability estimates on a common pooled estimate, so that additional
pseudocounts cannot masquerade as useful basis information. A measurement-balanced circuit
and an independent quantum-channel noise component test where that information disappears.

A second way to obtain information about disturbance is to acquire measurements at several
device-noise settings. Extrapolation can avoid an Eve-free reference run, but its precision
depends on setting separation and signal allocation. We connect this acquisition problem to
the reconciliation study through a common principle: extra information is useful only after
its sampling and modelling costs are included.

Our contributions are the following controlled evaluations.
\begin{enumerate}
\item We identify the complete-path mechanism behind the basis asymmetry and test its
attenuation when basis-independent channel noise replaces gate noise at fixed total QBER.
\item We measure the disclosure--FER tradeoff using nested sparse codes, independently
calibrated operating points, paired holdout frames, smoothing ablations, and a stronger
decoder. This establishes when basis information helps and when its finite-sample cost wins.
\item We compare calibrated and calibration-free disturbance estimators at a matched signal
budget, optimize acquisition design before simulation, and quantify drift sensitivity and
the ambiguity introduced by unscaled channel disturbance.
\end{enumerate}
The contribution is the interaction of circuit exposure, finite-sample inference, and
resource allocation. Position-dependent soft information and zero-noise extrapolation are
established tools; the depolarizing identity used to connect them to this implementation is
standard.

\section{Setting and Related Work}
\label{sec:setting}
BB84~\cite{bennett1984} retains matching-basis measurements and discloses a sample for parameter
estimation. We use uniformly random bits and bases, independent single-qubit signals, and an
authenticated noiseless classical channel. A reconciliation frame contains $n+S=1200$ sifted
bits, with $S=200$ randomly disclosed and $n=1000$ decoded. The evaluation conditions on this
sifted length. A separate diagnostic experiment counts every transmitted signal, including
signals discarded by sifting.

The Monte Carlo simulation samples each reconciliation bit's error from its exact
circuit-path probability. Diagnostic sifted counts follow $N\sim\mathrm{Binomial}(T,1/2)$
for $T$ transmitted signals, then error counts follow $E\sim\mathrm{Binomial}(N,Q)$.
Under the independent-shot model these distributions reproduce individual circuit
acquisitions; exact density-matrix checks validate the probabilities used by the sampler.

LDPC syndrome reconciliation and rate adaptation are established in
QKD~\cite{elkouss2011,blind2012,martinez2013}. Mesiti et al.~\cite{mesiti2010} combine quantum-
and public-channel metrics in sparse-graph decoding; puncturing and shortening produce
position-dependent priors~\cite{blind2012}. Origlia et al.~\cite{origlia2026} construct soft
metrics for reverse reconciliation in discrete-modulation continuous-variable QKD.
Basis-aware sifting and reconciliation~\cite{martinezmateo2024} exploit related protocol
structure. These works establish the value of soft information, while the measurement and
disclosure mechanisms differ from using an already-public BB84 basis label. Modern
reconciliation evaluations jointly report leakage and frame recovery
\cite{borisov2023,treeviriyanupab2024}.

Zero-noise extrapolation estimates idealized observables through additional acquisitions
\cite{temme2017,li2017,giurgicatiron2020,krebsbach2022}. Noisy-channel interception estimators
\cite{fiorini2024,fiorini2026} and detector-based separation of noise contributions
\cite{qi2007dualdet} provide relevant alternatives. We implement the published
interception-density formula alongside calibrated inversion, rather than treating measured
total QBER as an Eve-only estimator.

Our outcomes are reconciliation performance and model-based disturbance estimates.
Implementation security and finite-key extraction require source and detector assumptions,
an entropy bound, and privacy amplification
\cite{gllp2004,scarani2009,xu2020,fung2009,tomamichel2012,rennerkoenig2005}.
We therefore report reconciliation yield separately from a secure-key rate.

\section{Reliability Starts with the Complete Circuit}
\label{sec:mechanism}
\subsection{Receiver gates complete the exposure count}
Figure~\ref{fig:circuits} shows the actual preparation and matched-basis measurement paths.
PRIMARY prepares rectilinear zero without a gate, uses $X$ for rectilinear one, and uses $H$
or $XH$ for the diagonal states. Scheme~A adds an identity to every preparation, so every
path has nonzero exposure. Scheme~B uses two preparation gates for each state. Bob still
applies $H$ before diagonal measurement; B+ adds an identity before rectilinear measurement
to balance the complete paths. Each identity is retained as an explicit noise site.

% Identity operations are experimental noise carriers. Compiler removal would
% change the physical model; circuit exports retain all operations explicitly.
\begin{figure}[htbp]
\centering
\includegraphics[width=0.96\linewidth]{images/complete_circuits.pdf}
\caption{Complete source circuits, with gates ordered from left to right. The dotted boundary separates Alice's preparation from Bob's measurement rotation; the number at right counts noisy gates. All paths start in $|0\rangle$, and the indicated bit determines the preparation. Orange gates belong to Bob. B+ adds a noisy identity on his rectilinear path. Measurement itself is noiseless.}
\Description{Four encodings, each with two bit paths in each basis. PRIMARY has complete counts 0/1 and 2/3; A has 1/2 and 3/4; B has 2/2 and 3/3; B+ has 3/3 and 3/3.}
\label{fig:circuits}
\end{figure}

For a homogeneous depolarizing channel
$\mathcal D_p(\rho)=(1-p)\rho+pI/2$ after each gate, covariance under ideal gates gives
\begin{equation}
 q_k(p)=\frac{1-(1-p)^k}{2}.
 \label{eq:gate}
\end{equation}
The complete gate counts for the two bit values are $(0,1;2,3)$ in PRIMARY,
$(1,2;3,4)$ in A, $(2,2;3,3)$ in B, and $(3,3;3,3)$ in B+, where the semicolon separates
rectilinear and diagonal bases. The basis error probability is the average over its two
bit paths. In B, the ratio is $q_3(p)/q_2(p)$; at 10\% total target QBER it is 1.437.
The additional receiver Hadamard explains the residual asymmetry. Exact density-matrix
checks verify all paths and the two-hop intercept-resend model used below.

Scheme~A is a mechanism control rather than a proposed encoding optimization. Its identities
remove a literal zero-noise path while preserving exposure differences. B+ supplies the
corresponding negative control for basis information. Together they separate the source
of the effect from the special bare-state path in PRIMARY.

\subsection{A basis-independent channel attenuates the signal}
We insert an independent depolarizing link between preparation and reception, with bit-error
probability $c$ in either basis. Isotropy allows this channel to be commuted to the end of
the path for calculation. The resulting basis probabilities are
\begin{equation}
 p_z^{\mathrm{tot}}=c+(1-2c)p_z^{\mathrm{gate}},\qquad z\in\{r,d\}.
 \label{eq:channel}
\end{equation}
At fixed total QBER $Q$, the average gate contribution is $(Q-c)/(1-2c)$. Increasing $c$
therefore reduces the basis contrast, reaching zero at $c=Q$. This control isolates noise
composition from an increase in overall decoding difficulty.

\begin{figure}[htbp]
\centering
\includegraphics[width=0.94\linewidth]{images/gate_exposure.pdf}
\caption{Left: measured basis QBERs across the main holdout frames; circles indicate rectilinear and squares diagonal. Colors follow PRIMARY, A, B, and B+ as in the legend at right. Short horizontal marks give exact gate-model predictions. Wilson 95\% error bars are smaller than most markers because each point pools approximately 1.2 million bits. Right: predicted attenuation of the basis ratio when channel noise replaces gate noise at fixed 10\% total QBER.}
\Description{Basis error probabilities agree with gate-count predictions; independent channel noise drives each diagonal-to-rectilinear ratio to one.}
\label{fig:exposure}
\end{figure}

The error-conditioned Alice-bit proportion also follows from the same model,
$\Pr(X=1\mid\mathrm{error},z)=q_{k_{z1}}/(q_{k_{z0}}+q_{k_{z1}})$.
It is one in PRIMARY's rectilinear basis, approximately 0.657 in A at 10\% target QBER,
and one half in B and B+. These are channel-characterization quantities. The decoder
receives sampled basis probabilities, and receives no error positions or undisclosed bits.

\section{Estimating Reliability and Choosing a Code}
\label{sec:methods}
\subsection{A common centre removes the pseudocount confound}
Let $E_z$ errors be observed among $S_z$ disclosed bits in basis $z$, with
$E=E_r+E_d$ and $S=S_r+S_d$. The pooled Jeffreys estimate is
\begin{equation}
 \widehat p=\frac{E+1/2}{S+1}.
 \label{eq:pool}
\end{equation}
Applying Jeffreys smoothing independently in two bases adds a total of one pseudo-error,
whereas \eqref{eq:pool} adds one half. A direct pooled-versus-separate comparison thus
changes both the reliability contrast and the smoothing convention.

% Centre every class on the same pooled estimate. The positive-part weight is
% a plug-in shrinkage rule; no Bayes-optimality or universal dominance is claimed.
We instead use a pooled-centred positive-part shrinkage rule. With
$\widetilde p_z=E_z/S_z$, $\widetilde p=E/S$, and
$d=\widetilde p_d-\widetilde p_r$, set
\begin{align}
 v&=\widehat p(1-\widehat p)(S_r^{-1}+S_d^{-1}), &
 w&=\max\{0,1-v/d^2\}, \label{eq:weight}\\
 \widehat p_z^{\mathrm{sh}}&=\widehat p+w(\widetilde p_z-\widetilde p).
 &&\label{eq:shrink}
\end{align}
We set $w=0$ when $d=0$ or a class is absent. The sample-weighted average of the two
estimates remains $\widehat p$. The plug-in weight suppresses a contrast that is small
relative to its estimated variance. It is an empirical shrinkage heuristic, whose performance
is measured below. Ablations retain separate Jeffreys estimates and the coherent unshrunk
rule $w=1$ to distinguish centring from contrast shrinkage.

Bob estimates the error vector from the public syndrome difference and basis labels, using
$L_i=\log[(1-\widehat p_{z_i})/\widehat p_{z_i}]$.
Probabilities are clipped to $[10^{-6},0.499]$ for numerical stability, identically for all
decoders. Each frame has its own disjoint sample; the undisclosed frame is never used to
select its priors.

\subsection{A nested sparse family and a stronger decoder}
We construct a column-weight-three Tanner graph by progressive-edge-growth-style
breadth-first cycle avoidance and degree balancing~\cite{hu2005}. The initial matrix has
150 checks on 1000 variables. Repeatedly splitting the largest check into two disjoint
parts gives matrices with $m=150,175,\ldots,900$, preserving column weight three and
minimum check degree two. Every previous check is an XOR of later checks, so additional
rows add constraints within one nested family. The construction is fully specified by the
algorithm and matrix seed; it is a finite-length test family rather than an optimized
capacity-approaching QKD code.

Layered sum-product belief propagation (BP) uses at most 300 sweeps, stopping on exact
syndrome agreement. A stronger baseline adds full order-one ordered-statistics decoding
(OSD)~\cite{fossorier1995}. On BP nonconvergence or failed public verification, it selects
a least-reliable independent parity basis, re-encodes the BP hard decision, and tests each
single free-variable flip. Candidates are scored by the channel-prior LLRs. Both pooled
and basis-dependent methods receive the same iteration and search budgets.

Independent 64-bit Toeplitz verification determines whether Bob accepts a recovered frame.
Exact simulator correctness is recorded separately. Verification failure, rather than oracle
knowledge of the error pattern, triggers the OSD retry. The 64 public verification bits
are charged to the disclosure budget. Public random hash seeds contain no key information;
authentication cost is outside this reconciliation metric.

\subsection{Calibration precedes the independent evaluation}
For each encoding and target $Q\in\{0.02,0.06,0.10\}$, we generate 500 calibration frames.
Pooled and shrinkage BP policies each select the smallest tested $m$ for which the Wilson
95\% upper FER bound is at most 1\% and all higher tested row counts also pass. The initial
grids are $150{:}25{:}300$, $350{:}25{:}550$, and $500{:}25{:}800$, respectively.
If needed, the grid extends to 900. Choices are frozen before generating 2000 independent
holdout frames per cell. This is a conservative empirical operating-point rule; the
holdout data provide its separate performance assessment.

At each selected matrix, all four prior rules are paired on the same frames. We report both
BP and BP+OSD. The latter uses the BP-selected matrices, so this comparison tests decoder
strength at matched disclosure rather than optimizing an additional OSD rate policy.
Two additional matrix seeds, 29 and 47, use the same frozen row counts on fresh 1000-frame
cohorts at 10\% QBER. Fixed-code curves use 500 frames at each QBER from 1\% to 12\% in
A and B+, with $m=225,450,650$. Channel-composition controls have their own calibration
and 2000-frame holdout cohorts at $c=2.5,5,7.5,10\%$ and total $Q=10\%$.

\subsection{Report failure, disclosure, and uncertainty together}
With $I_{\mathrm{corr}}$ denoting exact recovery, we report
\begin{align}
 \mathrm{FER}&=1-\mathbb E[I_{\mathrm{corr}}], &
 Y&=\frac{n-m-64}{n+S}I_{\mathrm{corr}}, \label{eq:yield}\\
 f_{\mathrm{syn}}&=\frac{m}{n h_2(Q)}, &
 f_{\mathrm{EC}}&=\frac{m+64}{n h_2(Q)}. \label{eq:efficiency}
\end{align}
Here $h_2$ is binary entropy, and $Q$ in aggregate efficiency reporting is the measured
undisclosed-block QBER. The row count is charged conservatively even if checks are
dependent; algebraic ranks accompany the matrices. The algebraic code rate is
$R=1-\operatorname{rank}_{\mathrm{GF}(2)}(H)/n$; all tested matrices have full row rank.
$Y$ is a reconciliation yield per
sifted bit, including estimation and verification costs, before privacy amplification.

Absolute FERs include failure counts and Wilson 95\% intervals~\cite{wilson1927},
restricted to $[0,1]$ at numerical boundaries.
Same-matrix comparisons retain rescued and harmed pairs and use exact two-sided McNemar
tests. Holm correction is applied across the 12 main cells separately for each decoder
and matched-code comparison family. Paired bootstrap intervals use 10,000 resamples.
For operating-policy yield comparisons, intervals condition on the frozen calibration
choices; they do not include variation from repeating calibration.

\section{What Survives in a Low-FER Operating Regime?}
\label{sec:results}
\subsection{Code quality changes the relevant question}
Figure~\ref{fig:curves} shows the fixed-code transitions. Larger check counts move useful
operation to higher QBER, without mixing unrelated degree profiles into a rate ladder.
The calibrated operating points in Table~\ref{tab:operating} yield observed FER below 1\%
for both policies in all 12 main cells. Finite holdout counts still matter: a point estimate
below the target does not imply that every 95\% upper bound is below it.

\begin{figure}[htbp]
\centering
\includegraphics[width=0.94\linewidth]{images/fixed_code_curves.pdf}
\caption{Fixed-code BP curves, with 500 independent frames at each QBER. Solid lines use pooled priors and dashed lines use centred shrinkage. Each color is one matrix, reused across QBER. The dotted line marks 1\% FER; the vertical axis is linear below 1\% and logarithmic above it. Zero observed failures have Wilson 95\% upper bound 0.76\%. Full counts accompany the curves.}
\Description{FER-versus-QBER transitions for three fixed matrices in Scheme A and balanced control B+.}
\label{fig:curves}
\end{figure}

% Operating and matched-code tables refer to the same holdout frames. Their
% different questions are explicit, preventing nominal-QBER table collisions.
@OPERATING_TABLE@

The disclosure costs are substantial at short block length. At 10\% QBER in A, the
selected syndrome counts are @A10_MU@ and @A10_MA@ for pooled and shrinkage BP, respectively.
Their $f_{\mathrm{EC}}$ values are @A10_FU@ and @A10_FA@. At 2\%, conservative calibration
can cost more than the benefit of a small basis contrast. These efficiencies place the
results in a practical coding context and prevent a large failure ratio at an unsuitable
rate from being interpreted as a large protocol improvement.

\subsection{Separate decoding effects from selected-rate effects}
Table~\ref{tab:matched} compares the priors on the same shrinkage-selected matrix at 10\%
QBER, including the stronger decoder. Table~\ref{tab:operating} instead compares two
independently calibrated policies. At the operating point in A, their BP yield difference
is @A10_Y@ percentage points (95\% interval @A10_YCI@). The gain mainly reflects a smaller
selected syndrome, alongside residual frame failures. It is conditional on this calibration
realization. It should therefore be interpreted together with the same-matrix comparison,
the matrix replications, and the balanced control.

@MATCHED_TABLE@

Basis labels have the clearest value when the physical contrast is large. PRIMARY provides
the strongest contrast, while A removes its special noise-free preparation. B has only the
receiver-induced difference, and B+ has equal class probabilities. In B+, the sample can
still suggest an apparent contrast; shrinkage reduces its influence but cannot eliminate
finite-sample mistakes. The operating yield difference in B+ at 10\% is @B10_Y@ percentage
points. This control makes the estimation cost part of the result.

@ABLATION_TEXT@
The ablation counts in Appendix~\ref{app:ablation} show the size of each effect without
crediting additional smoothing as basis information. Replication across two further
matrices tests the persistence of the matched-rate effect at frozen disclosure; it also
shows why one calibrated rate choice is not a general optimum.

\subsection{Channel noise limits the value of the label}
Figure~\ref{fig:operating} combines operating-policy differences with a fixed-disclosure
channel-composition sensitivity analysis. The latter reuses the original gate-only
shrinkage-selected matrices, avoiding the failure floor introduced by adaptive code selection
in the channel cohorts. Its code choices therefore come from separate calibration rather
than optimization on the channel outcomes. The physical ratio approaches one as channel
noise replaces gate noise. At $c=10\%$, all errors come from the independent basis-symmetric link.
At that endpoint, shrinkage has 110 versus 98 failures in PRIMARY and 33 versus 27 in A
on the fixed codes. These losses expose the cost of fitting a contrast that has no
physical basis. The control demonstrates why gains from a gate-only simulator
must be interpreted through the noise composition of an intended implementation.

\begin{figure}[htbp]
\centering
\includegraphics[width=0.94\linewidth]{images/operating_points.pdf}
\caption{Paired BP comparisons. Left: operating-policy reconciliation-yield differences in the 12 main cells. Right: a fixed-disclosure sensitivity analysis at 10\% total QBER, retaining the gate-only calibrated shrinkage code ($m=625$ for PRIMARY and $650$ for A). The independent channel component increases on 2000 holdout frames per cell. Positive values indicate fewer shrinkage failures. Error bars are paired 95\% bootstrap intervals conditional on the selected matrices.}
\Description{Operating-policy yield differences show both gains and losses and their dependence on channel composition.}
\label{fig:operating}
\end{figure}

\section{Acquiring Information about Disturbance}
\label{sec:diagnostics}
\subsection{A target shared by all estimators}
Eve independently intercepts a fraction $r=\pEve$ of signals, measures in a random basis,
and resends that result. The ideal-device disturbance is $e_{\mathrm{Eve}}=r/4$. All
diagnostic estimators target this quantity under a fixed intercept-resend model. Every
diagnostic bit is disclosed, so these acquisitions produce no key.

Let $q_0(p)$ be PRIMARY's mean QBER without Eve and $q_I(p)$ the mean conditional on
interception, including noisy Alice--Eve and Eve--Bob paths. Then
\begin{equation}
 Q(p,r)=(1-r)q_0(p)+r q_I(p),\qquad
 q_0(p)=\frac14\sum_{k=0}^{3}q_k(p).
 \label{eq:attack}
\end{equation}
Appendix~\ref{app:diagnostic} gives the explicit two-hop expression. In particular,
$Q(0,r)=r/4$. Exact density-matrix calculations validate the model at the largest
noise parameter used in the acquisition design as well as at the nominal setting.

At an overall budget of 8000 transmitted signals, four-setting extrapolation uses
$f=(1,1.5,2,2.5)$ with 2000 signals each and scales $p$ to $fp$. Weighted linear fitting
estimates the intercept; a gate-model fit jointly estimates $(p,r)$ using \eqref{eq:attack}.
Estimates remain unclipped so signed bias and boundary dispersion remain visible.
The base signal noise is $p=0.03$, with 1000 replicates at each
$r\in\{0,0.1,0.2,0.3,0.4\}$.

Independent calibrated inversion acquires an Eve-free batch and an attacked batch at the
unscaled setting. With observed probabilities $\widehat q_{\mathrm{hw}}$ and $\widehat Q$,
we compare
\begin{align}
 \widehat e_{\mathrm{BSC}}&=\frac{\widehat Q-\widehat q_{\mathrm{hw}}}{1-2\widehat q_{\mathrm{hw}}},
 \label{eq:bsc}\\
 \widehat e_{\mathrm{gate}}&=\frac{\widehat Q-q_0(\widehat p)}{4[q_I(\widehat p)-q_0(\widehat p)]},
 \quad q_0(\widehat p)=\widehat q_{\mathrm{hw}}, \label{eq:calgate}\\
 \widehat e_{\mathrm F}&=\frac{\widehat Q-\widehat q_{\mathrm{hw}}}{4(1/4-\widehat q_{\mathrm{hw}}^2)}.
 \label{eq:fiorini}
\end{align}
Equation~\eqref{eq:fiorini} implements the interception-density formula in Fiorini et al.'s
public simulator~\cite{fiorini2026}, divided by four to use the shared disturbance target.
Its readout/channel-noise convention differs from our two-hop gate model. We additionally
simulate its corresponding algebraic law $Q=e_{\mathrm{ch}}+r(1/4-e_{\mathrm{ch}}^2)$ to check the comparator
under the assumptions it was designed for.

\subsection{Select acquisitions before measuring performance}
Closely spaced settings make the fitted gate noise and interception fraction strongly
correlated. We select a wider two-setting design by minimizing the worst Fisher-predicted
variance over $p\in\{0.02,0.03,0.04\}$ and the five attack fractions. Candidate high scales
are $2.5,4,6,8$, and low-scale budget shares range from 0.5 to 0.9. This analytic selection
precedes Monte Carlo evaluation. It selects $f=(1,8)$ and $(6400,1600)$ signals.
The calibrated allocation is selected on the same grid by a delta-method calculation,
giving 2400 Eve-free and 5600 attacked signals. Both designs pay the same total budget.

At $p=0.03$, the four-setting Fisher correlation of $(p,r)$ ranges from -0.949 to -0.958;
the two-setting design reduces its magnitude to 0.708--0.764. A correct nonlinear model
can therefore have little structural bias yet a large finite-count error when the
acquisition settings provide weak parameter separation. Improving the model alone does
not repair this variance cost.

@DIAGNOSTIC_TABLE@

Table~\ref{tab:diagnostic} separates mean absolute error (MAE), signed bias, and within-target
dispersion. The optimized two-setting gate fit has MAE @TWO_MAE@ percentage points, compared
with @FOUR_MAE@ for the original four-setting gate fit and @CAL_MAE@ for calibrated gate
inversion. The paired 95\% interval for the two-setting-minus-calibrated MAE difference is
@DIAG_DIFFCI@ percentage points. Wide scaling is effective in this homogeneous gate model;
its use in hardware requires a controllable range over which that model remains valid.

\subsection{Drift and unscaled noise answer different questions}
We vary the calibration-run parameter as $p_{\mathrm{cal}}=p_{\mathrm{signal}}(1+\delta)$,
including $\delta=\pm5,\pm10,\pm20\%$ and extending to $\pm50\%$. Signal noise remains
fixed. Figure~\ref{fig:diagnostic} shows calibrated performance against the same
calibration-free fit. The actual sampled comparison, rather than an assumed drift
advantage, determines whether a crossover occurs. @DRIFT_TEXT@

\begin{figure}[htbp]
\centering
\includegraphics[width=0.94\linewidth]{images/diagnostic_design.pdf}
\caption{Matched-budget diagnostic accuracy. Left: MAE versus interception fraction for the original and optimized acquisition designs. Right: pooled MAE under calibration drift; shading marks the requested $\pm20\%$ range. All methods use 8000 signals per replicate. The calibration-free horizontal reference assumes stationary gate noise during its own acquisitions. Pooled uncertainty is reported in Table~\ref{tab:diagnostic}; retained replicate counts support each curve.}
\Description{Wider settings improve gate-model extrapolation, while stale calibration increases inversion error.}
\label{fig:diagnostic}
\end{figure}

An independent receiver-channel component $c$ that is not scaled changes the zero-gate-noise
intercept to $c+(1-2c)r/4$. Thus an extrapolated intercept by itself cannot identify the
Eve-only contribution in the presence of unknown unscaled disturbance. Additional
simulations at $c=1\%$ and $3\%$ retain this confounding. This is distinct from calibration
drift: one concerns a stale reference, while the other concerns which physical noise
components the acquisition actually varies. Both matter when assigning meaning to an
estimated intercept.

\section{Interpretation and Scope}
\label{sec:discussion}
The reconciliation and diagnostic experiments expose two information costs. Public basis
labels are free to reuse, but the probabilities attached to them have finite estimation
error. Extra noise settings can remove the need for an Eve-free reference, but their
precision depends on the span and allocation of measurements. Circuit accounting explains
where the first source of information originates; experimental design determines whether
the second can be acquired efficiently.

For reconciliation, the useful result is a conditional disclosure advantage in a controlled
short-block study. The nested family and independent holdout remove the earlier dependence
on unrelated code profiles. The measured efficiencies also show the cost of this deliberately
conservative family. Larger blocks, optimized degree distributions, and rate policies
calibrated across repeated training cohorts are needed to establish a competitive practical
advantage. BP+OSD and the matrix replications check that the comparison is not solely a
weak-decoder effect; they do not substitute for that broader optimization.

The physical model uses homogeneous independent gate noise, ideal measurements, and an
isotropic channel. Gate-dependent errors, loss, detector imperfections, and correlations
can change the reliability classes. Physical identities represent explicit exposure sites;
a compiler that removes or reschedules them changes the experiment. Noise scaling likewise
assumes the attack and unscaled link remain stationary during acquisition. The channel
controls identify two concrete boundaries within these assumptions: attenuation of useful
basis contrast and ambiguity of the zero-gate-noise intercept.

Finally, reconciliation success and estimated disturbance have different security roles.
An adversary can adapt to announced settings, and a noise decomposition under a fixed
attack model is not a composable phase-error bound. A secure implementation must account
for all public leakage, correctness verification, source and detector assumptions, and
finite-key privacy amplification~\cite{shor2000,tomamichel2012,rennerkoenig2005}.
The measurements here support engineering comparisons within the stated model.
The 10\% cells are decoding stress tests: their positive reconciliation yield can be
consumed by subsequent privacy amplification.

\section{Conclusion}
Complete-path gate counts explain the basis asymmetry in this BB84 implementation, including
the receiver contribution that survives preparation balancing. With coherent probability
estimation and separately calibrated sparse codes, basis information has a conditional
disclosure benefit in a low-FER regime, alongside measurable estimation costs. Independent
channel noise narrows that opportunity. For disturbance estimation, wide, preselected noise
settings recover much of the precision lost by closely spaced extrapolation, while drift
and unscaled disturbance constrain the interpretation. These results turn implementation
details into testable conditions for when additional reliability information is useful.

\section*{Data and Code Availability}
The accompanying supplementary package contains the simulation and decoding code, all
parity-check matrices, explicit circuit exports, calibration and paired holdout records,
diagnostic counts, acquisition-design calculations, and scripts regenerating the tables,
figures, and manuscript. A frozen protocol, source provenance, dependency versions, and
file checksums document the inputs. Historical exploratory cohorts are archived separately
and do not contribute to the paper's numerical results.

\FloatBarrier
\bibliographystyle{ACM-Reference-Format}
\bibliography{references}
\FloatBarrier
\input{appendix}
\end{document}
